% Modernised reading — fonds Alexandre Grothendieck, shelfmark 97, pages 1-18
% (the whole folder, one batch).
%
% This is an INTERPRETATION, not a transcription. It re-reads the pages in
% today's notation and vocabulary, states what the manuscript leaves implicit,
% and moves everything the critical apparatus carried into footnotes. It is
% held to being mathematically correct as it stands; where the pages are loose
% or mistaken, the true statement is what appears here and a footnote says
% what the page has. In any dispute of substance the transcription governs:
% whoever cites Grothendieck must cite batch-01.fr.tex.
%
% Pass: Opus 5.5 (claude-opus-5-5), 2026-10-10 — first pass, unchecked against the pages by a human.
% The folder was read whole: one batch, pages 1-18, of which pages 2, 3, 5
% and 14 are not transcribed (a list of names; the abandoned opening of a
% letter to Atiyah on IHES letterhead; a blank; a blank verso). Nothing is
% said of pages 2 and 3 beyond what the transcription's header says.
%
% What the folder is. Under the catalogue title « Classes de Chern des
% représentations », three strands of his own and three sheets he wrote on:
%   p. 1        Serre's letter of 4 Nov. 1966 (sends a « Notice » and points
%               to a paper of Milnor on flat SL_2 bundles over a surface);
%               described, not restated.
%   p. 4        Chern classes of the regular representation of a finite
%               group, and H^{Nk}(G,Z)(p) != 0.
%   pp. 6-12    the flat bundle Sigma_0 on a compact hyperbolic surface X
%               given by a lift of the Fuchsian group to SL(2,R): Euler class
%               1-g, T_X = Sigma_0^{(x)2}, gamma_1^X = 2-2g; the lifting
%               obstruction in H^2(Gamma,Z/2); T_X + 2 flat.
%   p. 15       H^*(B_G,R) for G compact connected, via forms on the
%               simplicial manifold G^{p+1}/G, and Sym(T) = H^*(B_G,R).
%   p. 17       conditions « sans p-torsion » for a connected Lie group and
%               the implications between them (Hopf, Borel, Leray).
%   versos      pp. 7, 9, 11, 13: four sheets of a typescript on the Brauer
%               group of a fibration over a curve (typed pp. 23, 17, 19, 20);
%               p. 16: one sheet of a typed letter on the categories C(X,A);
%               p. 18: EGA IV typescript sheet IV-1102, (21.6.5)-(21.6.6).
%
% Notation fixed for the whole document. G = PSL(2,R) = SL(2,R)/{+-1} on
% pages 6-12 (he writes « GP(2,R) »), G~ = SL(2,R), S = PSO(2), S~ = SO(2);
% B_H the classifying space of a group H, as on the pages; gamma_1 the
% generator of H^2(B_S,Z), c_1 that of H^2(B_S~,Z); e( ) the Euler class;
% chi(X) for his EP(X). On page 17, B is a Borel subgroup and T a maximal
% torus. The typescript's Tate-Shafarevich group is written \mathrm{III}, as
% the transcription sets it.
%
% Substantive departures from the pages, each footnoted where it occurs:
% — p. 4: the product of the lambda in F_p^* is -1 (Wilson), not 1; so
%   c_{p-1}(r_G) = -xi_0^{p-1} and c_{N(p-1)}(N r_G) = (-1)^N xi^{N(p-1)}.
%   The conclusion is unaffected. The last lines drop the factor N; restored.
% — pp. 6-12: « GP(2,R) » read as PSL(2,R); the sign of gamma_1^X depends on
%   orientation conventions, fixed here. The argument by a covering of degree
%   d (p. 10) is stated as what it proves: compatibility, and reduction of
%   the formula for X to the formula for X'. Page 12 is almost illegible;
%   the reading of the extension class and of H^2(Gamma,Z/2) is ours.
%   Remark (4): T_X + 2 is already trivial as a real bundle; the content is
%   the flat structure it gets from Sigma_0 (x) Sigma_0.
% — p. 13 (typed 20): « la plus petite extension galoisienne ... de degré
%   nu » is replaced by the separable closure of k(y) in the function field
%   of the component, of degree nu; the two agree when it is Galois (always
%   for k(y) finite). « Pseudo-fini au sens de Serre » is Serre's
%   « quasi-fini ».
% — p. 16: the question « tout foncteur exact ... proviendrait d'un unique
%   morphisme » is false as posed (the zero functor, the shift); footnoted.
% — p. 17: (ii bis) => (ii) by Betti numbers needs the degrees of the
%   generators to be the rational ones; footnoted.
%
% Cross-references in the group « Géométrie et topologie combinatoire »:
% folder 88, pages 11-14 (the hyperbolic plane as homogeneous space of
% SL(2,R)/{+-1}, Fuchsian groups); folder 72, pages 107-112 (transcription
% only: Stiefel-Whitney classes and the Brauer class of a quadratic form,
% the same kind of mod 2 obstruction as on page 12 here).
%
% What the folder announces and never establishes: the computation the
% typescript begins at the foot of p. 13 (« on trouve ... »); the end of the
% EGA proof at the foot of p. 18; the reconstruction question of p. 16; the
% proofs of the degeneration and of the theorem of p. 15 (stated, not
% proved); Serre's « Notice » itself, which is not in the folder.
\documentclass[11pt,a4paper]{article}
\input{../preamble/grothendieck}

\folder{97}
\pages{1}{18}
\dating{[vers 1966]}
\watermark{Édition de démonstration\\interprétation personnelle de l'œuvre}
\foldertitle{Classes de Chern des représentations : lettre (1966), notes manuscrites (s.d.) — lecture modernisée du dossier entier}

\begin{document}

\section*{Résumé}

\begin{resume}
Une représentation d'un groupe --- une façon de le faire agir par matrices
--- fabrique des objets géométriques : des \emph{fibrés vectoriels}, familles
d'espaces vectoriels posés au-dessus de chaque point d'un espace, comme les
plans tangents posés en chaque point d'une surface. Un fibré peut être
« tordu » d'une façon qu'aucune déformation ne défait --- le ruban de Möbius
est un fibré en droites tordu au-dessus du cercle --- et ses \emph{classes
caractéristiques}, en particulier les classes de Chern, sont les entiers ou
les classes de cohomologie qui mesurent cette torsion. Ces feuillets
demandent ce que ces mesures valent sur les fibrés qui viennent d'une
représentation.

Deux cas sont traités. Pour un groupe fini, la représentation la plus
naturelle --- le groupe agissant sur lui-même par permutations --- a des
classes de Chern qu'on calcule exactement ; elles ne s'annulent jamais en
certains degrés, et cela suffit à montrer que la cohomologie d'un groupe fini
dont l'ordre est divisible par un nombre premier $p$ a de la $p$-torsion en
une infinité de degrés.

Le second cas est plus surprenant. Sur une surface compacte de genre
$g \geq 2$ --- une bouée à $g$ trous ---, la géométrie hyperbolique fournit
une représentation du groupe fondamental dans les matrices réelles $2 \times
2$ de déterminant $1$, et donc un fibré en plans sur la surface. Ce fibré est
\emph{plat} : on peut le transporter le long des chemins sans qu'il tourne,
comme si l'on n'avait jamais quitté un point. Pour des fibrés complexes
ordinaires, être plat force les classes de Chern à être négligeables ---
c'est ce que dit la théorie de Chern--Weil, qui calcule ces classes par la
courbure. Ici non : le fibré plat obtenu a une classe d'Euler égale à
$1 - g$, qui n'est pas nulle, et son « carré » est le fibré tangent de la
surface, de classe $2 - 2g$, la caractéristique d'Euler. L'analogie est celle
d'un escalier en colimaçon : localement chaque marche est plate, et pourtant
l'ensemble monte. La lettre de Serre qui ouvre le dossier renvoie à Milnor,
qui avait étudié en 1958 exactement ces fibrés plats ; la surface hyperbolique
réalise le cas extrême de son inégalité.

Autour de ces deux calculs, deux feuillets sur la cohomologie des espaces
classifiants des groupes de Lie, où ces classes vivent : à coefficients
réels, on la calcule par des formes différentielles invariantes, et elle est
une algèbre de polynômes ; modulo $p$, elle l'est encore exactement quand le
groupe n'a pas de $p$-torsion. Le reste du dossier est le papier sur lequel
il écrivait : des feuillets d'un tapuscrit sur le groupe de Brauer d'une
surface fibrée sur une courbe, une page d'EGA IV, et une page d'une lettre
dactylographiée. Cette dernière vaut d'être lue pour elle-même : elle
propose de regarder un espace non pas directement, mais à travers la
catégorie de tous les coefficients qu'il peut porter, et se demande si
cette catégorie suffit à le reconstituer --- « je crois que quelqu'un devrait
bien un jour regarder ces choses systématiquement ». C'est un programme
qu'il reprendra vingt ans plus tard dans \emph{À la poursuite des champs}.

\keywords{Chern classes of representations, regular representation, cohomology of finite groups, Fuchsian group, flat bundle, Euler class, Milnor-Wood inequality, spin structure, Adams operation, classifying space, simplicial de Rham complex, Koszul duality, torsion in Lie groups, Borel transgression theorem, Brauer group of a fibred surface, Tate-Shafarevich group, local systems and homotopy type, Weil divisor of a Cartier divisor}
\end{resume}

\section*{Le fil du dossier, et les conventions}

Le dossier n'a pas une argumentation, mais trois brins de même sujet, posés
sur des feuillets de remploi.

\begin{enumerate}
\item[1.] \emph{La lettre de Serre} (page 1), datée du 4 novembre 1966.
\item[2.] \emph{Classes de Chern de la représentation régulière} (page 4) :
groupes finis.
\item[3.] \emph{Le fibré plat d'une surface hyperbolique} (pages 6, 8, 10 et
12, quatre feuillets numérotés par lui 1, --, 2, 3) : groupes discrets dans
$\mathrm{PSL}(2,\mathbb{R})$.
\item[4.] \emph{Les espaces classifiants des groupes de Lie} (pages 15 et 17) :
coefficients réels, puis la $p$-torsion.
\end{enumerate}

Les versos des feuillets manuscrits portent d'autres textes, qui n'ont pas de
rapport avec les classes de Chern et qu'on lit à leur place dans l'ordre des
pages : un tapuscrit sur le groupe de Brauer d'une fibration (pages 7, 9, 11
et 13, qu'on lit dans l'ordre de leur pagination dactylographiée, 17, 19, 20,
23), une page de lettre dactylographiée sur les catégories $C(X,A)$
(page 16), une page du tapuscrit d'EGA IV (page 18). Les pages 2 (une liste
de noms), 3 (le début abandonné d'une lettre à Atiyah sur papier de l'IHÉS),
5 (blanche) et 14 (verso blanc) ne sont pas transcrites et ne sont pas lues
ici.

\emph{Conventions.} Pour un groupe topologique $H$, $B_{H}$ est son espace
classifiant, comme sur les pages. Aux pages 6 à 12, $G = \mathrm{PSL}(2,
\mathbb{R}) = \mathrm{SL}(2,\mathbb{R})/\{\pm 1\}$ opère sur le disque $D$
(ou le demi-plan de Poincaré) avec stabilisateur $S = \mathrm{PSO}(2)$ ;
$\widetilde{G} = \mathrm{SL}(2,\mathbb{R})$ et $\widetilde{S} =
\mathrm{SO}(2)$ sont leurs revêtements \emph{doubles} --- et non universels ---
et $\Gamma \subset G$ est un sous-groupe discret sans torsion, cocompact, de
quotient $X = D/\Gamma$ de genre $g \geq 2$. La caractéristique d'Euler
qu'il note $EP(X)$ est notée $\chi(X) = 2 - 2g$, et $e(\,\cdot\,)$ est la
classe d'Euler d'un fibré réel orienté de rang 2, identifiée à $c_{1}$ du
fibré en droites complexes associé. On identifie $H^{2}(X,\mathbb{Z})$ à
$\mathbb{Z}$ par la classe fondamentale de l'orientation complexe. À la
page 17, $B$ est un sous-groupe de Borel et $T \subset B$ un tore maximal.

\pagerange{1}{1}
\section*{La lettre de Serre (page 1)}

La page 1 est une lettre de J.-P.\ Serre à Grothendieck, sur papier à en-tête
du Collège de France, datée du 4 novembre 1966. Serre n'a pas regardé le
problème que Grothendieck lui avait posé, mais lui envoie une « Notice » qui
démontre la même chose, et qui renvoie à un papier de Milnor sur les fibrés
plats de groupe $\mathrm{SL}_{2}$ au-dessus d'une surface.\footnote{La lettre
est transcrite entière dans \texttt{batch-01.fr.tex} ; on ne la redonne pas
ici, n'étant pas de Grothendieck. « La même chose » est une lecture
incertaine. Ni le problème posé ni la « Notice » ne sont dans le dossier.
Le papier de Milnor est très vraisemblablement « On the existence of a
connection with curvature zero », \emph{Comment.\ Math.\ Helv.}\ 32 (1958),
qui montre qu'un fibré réel orienté de rang 2 sur une surface de genre $g
\geq 1$ admet une structure plate de groupe $\mathrm{GL}^{+}(2,\mathbb{R})$
si et seulement si sa classe d'Euler $e$ vérifie $|e| \leq g - 1$ ; l'inégalité
a été étendue aux fibrés plats en cercles par J.\ Wood (1971). L'identification
est la nôtre : la lettre ne donne pas de titre.} Le seul énoncé mathématique
de la lettre nomme le sujet des feuillets 6 à 12 ; c'est par eux que le
dossier se rattache à 1966.

\pagerange{4}{4}
\section*{Classes de Chern de la représentation régulière (page 4)}

Le titre est de sa main : « Classe de Chern de la représentation régulière et
ses multiples ».

\emph{Le groupe cyclique.} Soit $G$ cyclique d'ordre $n$ et $\check{G} =
\mathrm{Hom}(G, \mathbb{Q}/\mathbb{Z})$ son dual. Le cobord de la suite
$0 \to \mathbb{Z} \to \mathbb{Q} \to \mathbb{Q}/\mathbb{Z} \to 0$ donne
$H^{2}(G,\mathbb{Z}) \simeq H^{1}(G,\mathbb{Q}/\mathbb{Z}) =
\check{G}$,\footnote{Le « $\mathrm{Hom}$ » de la page est une lecture
incertaine : le signe ressemble à un $H$ surmonté d'une étoile. L'énoncé est
juste sous l'une ou l'autre lecture.} et l'on a, en degrés strictement
positifs,
\[
  H^{*}(G,\mathbb{Z})^{+} \simeq \mathrm{Sym}_{\mathbb{Z}/n\mathbb{Z}}(\check{G})^{+}
  \simeq (\mathbb{Z}/n\mathbb{Z})[\xi]^{+},
\]
où $\xi$ est un générateur de $H^{2}(G,\mathbb{Z})$, de degré 2 ; les degrés
impairs sont nuls. La représentation régulière complexe $r_{G}$ est la somme
des caractères $\chi \in \check{G}$, et la classe de Chern d'un caractère est
son image dans $H^{2}$ ; la formule de Whitney donne donc, pour un générateur
$\xi_{0}$ de $\check{G}$,
\[
  c(r_{G}) = \prod_{\chi \in \check{G}} (1 + \chi)
  = \prod_{\lambda \in \mathbb{Z}/n\mathbb{Z}} (1 + \lambda \xi_{0}) .
\]

\emph{Le cas $n = p$ premier.} Le facteur $\lambda = 0$ vaut 1, d'où
$c_{p}(r_{G}) = 0$, et
\[
  c_{p-1}(r_{G}) = \Bigl(\prod_{\lambda \in \mathbb{F}_{p}^{*}} \lambda\Bigr)
  \xi_{0}^{p-1} = -\,\xi_{0}^{p-1},
\]
indépendamment du choix de $\xi_{0}$, puisque $\lambda^{p-1} = 1$ pour
$\lambda \in \mathbb{F}_{p}^{*}$.\footnote{La page écrit $c_{p-1}(r_{G}) =
\xi_{0}^{p-1}$, le produit des éléments de $\mathbb{F}_{p}^{*}$ étant compté
pour 1. Par le théorème de Wilson il vaut $(p-1)! \equiv -1$. Le signe ne
change rien à ce qui suit, qui ne demande que la non-nullité.} Comme $c(N
r_{G}) = c(r_{G})^{N}$,
\[
  c_{i}(N r_{G}) = 0 \quad \text{pour } i > N(p-1), \qquad
  c_{N(p-1)}(N r_{G}) = (-1)^{N} \xi_{0}^{N(p-1)} \neq 0 ,
\]
cette dernière classe engendrant $H^{2N(p-1)}(G,\mathbb{Z}) =
\mathbb{Z}/p\mathbb{Z}$.

\emph{Le cas général.} Soit $G$ un groupe fini d'ordre $n$ divisible par $p$,
$H \subset G$ un sous-groupe d'ordre $p$ et $n' = n/p$. La restriction de
$r_{G}$ à $H$ est $n' r_{H}$, donc $(N r_{G})|_{H} = N n' r_{H}$ et
\[
  c_{N n'(p-1)}(N r_{G})|_{H} = c_{N n'(p-1)}(N n' r_{H}) \neq 0 .
\]
A fortiori $c_{N n'(p-1)}(N r_{G}) \neq 0$.\footnote{La page écrit ici
$c_{i}(r_{G})$ et $c_{N n'(p-1)}(r_{G})$, sans le facteur $N$ des lignes
précédentes, et une fois un $n$ minuscule qui est vraisemblablement ce $N$ ;
on le rétablit, l'énoncé sans $N$ étant faux pour $N > 1$ par raison de
degré ($c_{i}(r_{G}) = 0$ pour $i > n$).} Comme $H^{*}(H,\mathbb{Z})$ est de
$p$-torsion en degrés positifs, la restriction tue la partie première à $p$ ;
la composante $p$-primaire de cette classe est donc non nulle :
\[
  \boxed{H^{Nk}(G,\mathbb{Z})(p) \neq 0 \quad \text{pour tout entier } N \geq 1,
  \qquad k = 2n'(p-1).}
\]
En particulier la cohomologie entière d'un groupe fini a de la $p$-torsion en
une infinité de degrés pour tout premier $p$ divisant son ordre, et ces
classes sont des classes de Chern de représentations.\footnote{Les classes de
Chern des représentations d'un groupe fini avaient été introduites et
étudiées par Atiyah, « Characters and cohomology of finite groups »,
\emph{Publ.\ Math.\ IHÉS} 9 (1961) ; la page 3 du dossier porte le début
abandonné d'une lettre à Atiyah, dont rien n'est lisible au-delà de
l'en-tête et de l'adresse. Que la cohomologie d'un groupe fini ait de la
$p$-torsion en une infinité de degrés est classique ; la page ne cite
personne. Le rang $k$ obtenu n'est qu'une borne commode : pour $G$ cyclique
d'ordre $p$ il y a de la $p$-torsion en tout degré pair.}

\pagerange{6}{12}
\section*{Le fibré plat d'une surface hyperbolique (pages 6, 8, 10 et 12)}

Quatre feuillets à l'encre noire, numérotés par lui 1, --, 2, 3, écrits au dos
des feuillets du tapuscrit qui les suivent dans le dossier. Les pages 8 et 12
sont de l'écriture la plus rapide du dossier, et une grande partie de leur
prose ne se lit pas ; les formules se lisent. Ce qui suit est la mise au net
d'un argument dont les formules sont sur la page ; les phrases qui les
relient sont, en partie, les nôtres, et les notes disent où.

\pagerange{6}{6}
\subsection*{La situation, et le dictionnaire (page 6)}

Le disque $D$ est l'espace homogène $G/S = \widetilde{G}/\widetilde{S}$,
$S = G_{x}$ et $\widetilde{S} = \widetilde{G}_{x}$ étant les stabilisateurs
d'un point $x$. Le groupe $\Gamma$ se relève isomorphiquement en
$\widetilde{\Gamma} \subset \widetilde{G}$ (on le verra page 12), et l'on a le
diagramme
\begin{tikzcd}[column sep=small, row sep=small]
\widetilde{S} = \mathrm{SO}(2) \arrow[rr] \arrow[d, hook] & & S = \mathrm{PSO}(2) \arrow[d, hook] \\
\widetilde{G} = \mathrm{SL}(2,\mathbb{R}) \arrow[rr] & &
G = \mathrm{PSL}(2,\mathbb{R}) \\
\widetilde{\Gamma} \arrow[rr, "\sim"] \arrow[u, hook] & & \Gamma \arrow[u, hook]
\end{tikzcd}
dont les flèches horizontales sont de degré 2.\footnote{La page note $G =
\mathrm{GP}(2,\mathbb{R}) = \widetilde{G}/(+1,-1)$ ; « GP » est lu tel
qu'écrit, lecture incertaine, et ce qu'il désigne est sans ambiguïté
$\mathrm{PSL}(2,\mathbb{R})$. Le diagramme de la page place $\widetilde{S}
\to S$ au-dessus de $\widetilde{G} \to G$ sans flèches verticales, et des
traits sans pointe lisible entre $\widetilde{G}$, $G$ et $\widetilde{\Gamma}$,
$\Gamma$ ; on les rend par les inclusions. C'est le même plan hyperbolique,
espace homogène de $\mathrm{SL}(2,\mathbb{R})/\{\pm 1\}$, que le dossier 88
décrit aux pages 11 à 13 comme l'une des trois surfaces standard, et dont il
range les pavages réguliers à la page 14.}

\emph{Le dictionnaire.} Toute représentation linéaire réelle $W$ de
$\widetilde{S}$ définit un fibré vectoriel $\widetilde{G}$-équivariant
$\widetilde{G} \times^{\widetilde{S}} W$ sur $D$, donc, par restriction de
l'action à $\widetilde{\Gamma}$ et passage au quotient, un fibré vectoriel
réel $\mathcal{F}(W)$ sur $X$. La correspondance $W \mapsto \mathcal{F}(W)$
respecte sommes, produits tensoriels, puissances extérieures et dualité :
c'est un homomorphisme d'anneaux $\lambda$ de $RO(\widetilde{S})$ dans
$KO(X)$.\footnote{La page dit qu'elle « respecte les opérations tensorielles,
l'opération de conjugaison » ; la suite est illisible. Le nom d'anneau
$\lambda$ et la notation $\mathcal{F}$ sont les nôtres. En marge, à l'encre
brune, un mot biffé lu « Lusin », sans rapport visible avec le texte.} Deux
exemples :

\begin{itemize}
\item[--] la représentation de $\widetilde{S}$ sur $T_{x}D$, qui se factorise
par $S$, donne le fibré tangent $T_{X}$ ;
\item[--] si $W$ est la restriction d'une représentation $V$ de $\widetilde{G}$
tout entier, $\widetilde{G} \times^{\widetilde{S}} V \simeq D \times V$ avec
l'action diagonale $\gamma \cdot (z, v) = (\gamma z, \gamma v)$, et
$\mathcal{F}(V)$ est le fibré \emph{plat} associé à la représentation
$\widetilde{\Gamma} \to \widetilde{G} \to \mathrm{GL}(V)$ du groupe
fondamental.
\end{itemize}

Soit $V = \mathbb{R}^{2}$ la représentation standard de $\mathrm{SL}(2,
\mathbb{R})$, et $\Sigma_{0} = \mathcal{F}(V)$ le fibré plat de rang 2 qu'elle
définit. Restreinte à $\widetilde{S} = \mathrm{SO}(2) = \{e^{i\theta}\}$, $V$
est la rotation d'angle $\theta$, de caractère complexifié $e_{1} + e_{-1}$ ;
$T_{x}D$ est la rotation d'angle $2\theta$, de caractère $e_{2} + e_{-2}$.
D'où, dans l'anneau des représentations,
\[
  [T_{x}D] = \psi^{2}[V] = [V]^{2} - 2,
\]
où $\psi^{2}$ est la deuxième opération d'Adams ; et, par le dictionnaire,
\[
  \Sigma_{0} \otimes_{\mathbb{R}} \Sigma_{0} \simeq T_{X} \oplus
  2 \cdot \mathbf{1}_{\mathbb{R}} .
\]
\footnote{La page écrit $X = X_{0}^{(2)}$ et $X = X_{0}^{2} - 2$, avec
$X_{0}$, $X$ pour les représentations $V$ et $T_{x}D$ de $\widetilde{S}$, et
leurs caractères $e_{1} + e_{-1}$, $e_{2} + e_{-2}$. Lire $X_{0}^{(2)}$ comme
l'opération d'Adams $\psi^{2}$ est notre interprétation ; elle est conforme
à la formule $\psi^{2}(x) = x^{2} - 2\lambda^{2}(x)$, $\lambda^{2}V =
\det V$ étant trivial. La lettre $X$ désigne aussi la surface ; on ne la
reprend pas pour la représentation. L'isomorphisme de fibrés est écrit dans
la marge de la page 8, verticalement.} En termes de représentations de
$\mathrm{SL}(2,\mathbb{R})$, $V \otimes V = \mathrm{Sym}^{2} V \oplus
\Lambda^{2} V$, où $\Lambda^{2}V$ est triviale et $\mathrm{Sym}^{2}V$ est la
représentation adjointe, de restriction $e_{2} + e_{-2} + 1$ à
$\widetilde{S}$.

\emph{Réel, non complexe.} Les fibrés réels orientés de rang 2 sont aussi des
fibrés en droites complexes, $\mathrm{SO}(2) = \mathrm{U}(1)$ ; à ce titre
$T_{X}$ est le fibré tangent holomorphe et $\Sigma_{0}$ un fibré en droites
complexes $L_{0}$, avec
\[
  T_{X} \simeq L_{0}^{\otimes 2} \quad \text{(fibrés en droites complexes)},
\]
puisque le caractère $e_{2}$ est le carré de $e_{1}$. Mais la structure plate
de $\Sigma_{0}$ est réelle et non complexe : $\mathrm{SL}(2,\mathbb{R})$ ne
préserve aucune structure complexe de $\mathbb{R}^{2}$, et c'est ce qui va
permettre à $\Sigma_{0}$ d'avoir une classe caractéristique non
nulle.\footnote{La page insiste : « représentation réelle (pas complexe !) »
(page 8) et, page 6, « rep.\ complexe --- mais pas nécessairement \ldots »,
le dernier mot illisible. Un fibré en droites complexes muni d'une connexion
plate \emph{complexe} a une classe de Chern de torsion, par Chern--Weil ;
une connexion plate de groupe $\mathrm{SL}(2,\mathbb{R})$ ne se prête pas à
ce calcul, la classe d'Euler n'étant pas un polynôme de Chern--Weil pour le
groupe non compact $\mathrm{SL}(2,\mathbb{R})$.}

\pagerange{8}{8}
\subsection*{Le calcul : $e(\Sigma_{0}) = 1 - g$ (page 8)}

Écrivons $\gamma_{1}$ pour la classe de $T_{X}$ et $c_{1}$ pour celle de
$\Sigma_{0}$. De $T_{X} \simeq L_{0}^{\otimes 2}$ on tire
\[
  \gamma_{1} = 2 c_{1}.
\]
Or la classe d'Euler du fibré tangent est la caractéristique d'Euler,
$\gamma_{1} = e(T_{X}) = \chi(X) = 2(1-g)$, d'où
\[
  \boxed{c_{1} = e(\Sigma_{0}) = 1 - g \neq 0.}
\]
\footnote{La page note d'abord $\widetilde{c}_{1}$ et $c_{1}$ les classes de
$\Sigma = T_{X}$ et de $\Sigma_{0}$, puis écrit $\gamma_{1}$ pour la
première, qui est la notation du feuillet 10 ; on la garde. Les deux fibrés
y sont dits « inversibles », d'abord « réels », corrigé en « complexes ».
Le reste du feuillet --- sept lignes --- ne se lit que par mots isolés, parmi
lesquels « sur $D$ », « de $\Gamma$ », « de $\widetilde{S}$ », « les fibrés
$\Sigma_{0}$ », « $T_{X}$ », « d'une connexion intégrable » : il décrivait,
semble-t-il, la construction de $\Sigma_{0}$ par passage au quotient d'un
fibré sur $D$, et la conséquence qu'on lit en marge et à la page 12.}

Ainsi $\Sigma_{0}$ est un fibré plat de rang 2 dont la classe d'Euler n'est
pas nulle : c'est le cas d'égalité $|e| = g - 1$ de l'inégalité de Milnor, et
c'est ce qui rend la lettre de Serre pertinente.\footnote{Le rapprochement
avec Milnor et la lettre de Serre est le nôtre ; la page ne cite personne.
Topologiquement $L_{0}$ est une racine carrée du fibré tangent, c'est-à-dire
l'inverse d'une caractéristique thêta $K_{X}^{1/2}$ ; les relèvements de
$\Gamma$ à $\mathrm{SL}(2,\mathbb{R})$ correspondent aux $2^{2g}$ structures
spin de $X$. Le fibré tangent lui-même, de classe $|2 - 2g| = 2g - 2 \geq
g$, n'admet pas de structure plate (Milnor). Que la classe d'Euler d'une
représentation de $\pi_{1}(X)$ dans $\mathrm{PSL}(2,\mathbb{R})$ ne vaille
$\pm(2 - 2g)$ que pour les représentations fuchsiennes est un théorème de
Goldman (1980), bien postérieur.}

\emph{En marge.} Comme $\Sigma_{0} \otimes \Sigma_{0}$ est plat, $T_{X}
\oplus 2 \cdot \mathbf{1}_{\mathbb{R}}$ admet une connexion
intégrable ; on retrouve l'énoncé à la page 12.

\pagerange{10}{10}
\subsection*{La même chose dans les espaces classifiants (page 10)}

L'inclusion $\Gamma \subset G$ donne, puisque $D$ est contractile, une
application $X = B_{\Gamma} \to B_{G}$, et les inclusions des compacts
maximaux sont des équivalences d'homotopie $B_{S} \to B_{G}$, $B_{\widetilde{S}}
\to B_{\widetilde{G}}$. En cohomologie :
\begin{tikzcd}
H^{*}(B_{G},\mathbb{Z}) \arrow[r, "\sim"] \arrow[d] \arrow[rd] &
H^{*}(B_{S},\mathbb{Z}) = \mathbb{Z}[\gamma_{1}] \arrow[d, "\gamma_{1} \mapsto 2c_{1}"] \\
H^{*}(B_{\Gamma},\mathbb{Z}) = H^{*}(X,\mathbb{Z}) &
H^{*}(B_{\widetilde{S}},\mathbb{Z}) = \mathbb{Z}[c_{1}] \simeq H^{*}(B_{\widetilde{G}},\mathbb{Z})
\end{tikzcd}
où $\gamma_{1}$ est la première classe de Chern du caractère de $S$ sur
$T_{x}D$ et $c_{1}$ celle de la représentation standard de $\widetilde{S} =
\mathrm{U}(1)$ ; l'image de $\gamma_{1}$ est $2c_{1}$ parce que $\widetilde{S}
\to S$ est l'élévation au carré.\footnote{Le diagramme de la page dispose
$H^{*}(B_{G})$ en haut, $H^{*}(B_{\widetilde{G}})$ et $H^{*}(B_{S})$ en
dessous, $H^{*}(B_{\widetilde{S}})$ en bas à gauche, et une flèche ondulée de
$\mathbb{Z}[\gamma_{1}]$ vers $\gamma_{1}^{X}$ ; on l'a redressé en carré, la
flèche diagonale étant la composée $H^{*}(B_{G}) \to H^{*}(B_{\widetilde{G}})$.
Le « Car » qui précède $\gamma_{1} = 2c_{1}$ est une lecture incertaine.}
Notons $\gamma_{1}^{X} \in H^{2}(X,\mathbb{Z}) \simeq \mathbb{Z}$ l'image de
$\gamma_{1}$. Comme le $G$-fibré principal plat de $X$, réduit à $S$, a pour
fibré associé $T_{X}$, $\gamma_{1}^{X} = e(T_{X})$.

Le relèvement $\widetilde{\Gamma}$ donne un relèvement $X \to
B_{\widetilde{G}}$ de $X \to B_{G}$ :
\begin{tikzcd}
& B_{\widetilde{G}} \arrow[d] \\
X = B_{\Gamma} \arrow[r] \arrow[ru, dashed] & B_{G}
\end{tikzcd}
et l'image $c_{1}^{X}$ de $c_{1}$ est la classe d'Euler du fibré de rang 2
qu'il classifie, qui est $\Sigma_{0}$. D'où
\[
  \tfrac{1}{2}\gamma_{1}^{X} = c_{1}^{X} = e(\Sigma_{0}) = 1 - g, \qquad
  \boxed{\gamma_{1}^{X} = 2 - 2g = \chi(X) \neq 0.}
\]
\footnote{Les signes dépendent de l'orientation de $X$ et du choix du
générateur $\gamma_{1}$ ; on les fixe par la structure complexe de $D$,
pour laquelle $\gamma_{1}^{X}$ est la classe de Chern du fibré tangent
holomorphe. La page écrit « À calculer $\gamma_{1}^{X}$ ?? », puis, après
quelques mots illisibles, donne le résultat ; la phrase qui introduit $c_{1}$
comme « classe de Chern $c_{1}^{\mathbb{R}}$ d'un fibré vectoriel réel
$\Sigma$ sur $B_{\widetilde{G}}$, de rang 2 » est en partie illisible et en
partie incertaine.}

\emph{Le contrôle par un revêtement.} Soit $\Gamma' \subset \Gamma$ un
sous-groupe d'indice fini $d$ --- il en existe, $\Gamma$ étant résiduellement
fini ---, et $p : X' = B_{\Gamma'} \to X$ le revêtement étale de degré $d$
correspondant. La classe $\gamma_{1}^{X'}$ est $p^{*}\gamma_{1}^{X}$, et
$p^{*}$ est la multiplication par $d$ sur $H^{2}(\,\cdot\,,\mathbb{Z}) \simeq
\mathbb{Z}$ ; d'autre part $\chi(X') = d\,\chi(X)$. Les deux membres de
$\gamma_{1}^{X} = \chi(X)$ se multiplient donc par $d$ quand on passe à $X'$ :
la formule est compatible aux revêtements, et, $H^{2}(X,\mathbb{Z})$ étant
sans torsion, il suffit de l'établir pour un seul revêtement fini de $X$
pour l'avoir sur $X$.\footnote{La page conclut « on conclut $\gamma_{1}^{X} =
EP(X)$. OK. » ; les mots qui relient les formules sont en partie devinés par
la transcription, en particulier « comme $\gamma_{1}^{X'} = EP(X')$ », qui,
lu littéralement, présupposerait ce qu'on veut établir. On énonce ce que
l'argument prouve.}

\pagerange{12}{12}
\subsection*{La classe comme classe d'extension ; le relèvement (page 12)}

\emph{La classe d'extension.} Le revêtement universel $\widehat{G}$ de $G =
\mathrm{PSL}(2,\mathbb{R})$ est une extension centrale de $G$ par $\pi_{1}(G)
= \mathbb{Z}$ ; son image réciproque sur $\Gamma$ est une extension centrale
\[
  0 \to \mathbb{Z} \to \widehat{\Gamma} \to \Gamma \to 1,
\]
dont la classe vit dans $H^{2}(\Gamma,\mathbb{Z}) \simeq H^{2}(X,\mathbb{Z})$ :
c'est, au signe près, $\gamma_{1}^{X}$. Plus généralement, pour un groupe de
Lie connexe $G$ de groupe fondamental $H$, la fibration $B_{H} \to
B_{\widehat{G}} \to B_{G}$ et l'isomorphisme $H^{2}(B_{G}, H) \simeq
\mathrm{Hom}(\pi_{2}B_{G}, H) = \mathrm{Hom}(H,H)$ définissent une classe
canonique, image de l'identité, qui classifie le revêtement universel ; pour
$G = S$ un cercle, $H = \mathbb{Z}$, c'est la classe de Chern $\gamma_{1}$
de $H^{2}(B_{S},\mathbb{Z})$, et l'image réciproque à $\Gamma$ de la classe
canonique est la classe de l'extension.

\emph{Le relèvement.} De même le revêtement double $\widetilde{G} =
\mathrm{SL}(2,\mathbb{R}) \to G$ définit une classe dans $H^{2}(\Gamma,
\mathbb{Z}/2\mathbb{Z}) = \mathbb{Z}/2\mathbb{Z}$, qui est l'obstruction à
relever $\Gamma$ dans $\mathrm{SL}(2,\mathbb{R})$, et qui est la réduction
modulo 2 de $\gamma_{1}^{X}$. Comme $\gamma_{1}^{X} = \chi(X) = 2(1-g)$ est
pair, l'obstruction est nulle : $\Gamma$ se relève, ce qu'on a utilisé page
6.\footnote{La page 12 ne se lit presque que par ses formules :
« Le rev.\ universel de $G$ est de noyau $\mathbb{Z}$ », « extension de
$\Gamma$ par $\mathbb{Z}$ », $H^{2}(\Gamma,\mathbb{Z}) \simeq
H^{2}(X,\mathbb{Z})$, « que c'est $\gamma_{1}$ », $H^{2}(B_{G}, H)$,
$B_{H} \to B_{E} \to B_{G}$, « Chern », $H^{2}(\Gamma,\mathbb{Z}/2\mathbb{Z})
= \mathbb{Z}/2\mathbb{Z}$, « $\widetilde{G} \to G$ », « Comme $EP(X) =
2(1-g)$ ». L'articulation qu'on en donne est la nôtre ; elle est conforme à
ces fragments et elle est juste. La page note $E$ le revêtement universel,
qu'on note $\widehat{G}$ pour ne pas le confondre avec $\widetilde{G}$, qui
n'en est qu'un quotient. Le dossier 72, pages 107 à 112, calcule une
obstruction de même nature, en caractéristique arithmétique : la classe de
Brauer qui dit si un fibré en droites projectives vient d'un fibré vectoriel
de rang 2.}

\emph{Remarque (4).} Dans la situation (1) --- $X$ de genre $g \geq 2$, de
fibré tangent $T_{X}$ ---, le fibré vectoriel réel $T_{X} \oplus 2 \cdot
\mathbf{1}_{\mathbb{R}}$ sur $X$ est muni d'une connexion intégrable, celle
de $\Sigma_{0} \otimes \Sigma_{0}$, de groupe d'holonomie l'image de $\Gamma$
par $\mathrm{SL}(2,\mathbb{R}) \to \mathrm{GL}(V \otimes V)$.\footnote{En tant
que fibré réel, $T_{X} \oplus 2 \cdot \mathbf{1}_{\mathbb{R}}$ est déjà
trivial --- une surface orientable se plonge dans $\mathbb{R}^{3}$, d'où
$T_{X} \oplus \mathbf{1} \simeq \mathbf{3}$ --- et admet donc la connexion
triviale. Ce que la remarque apporte est une structure plate non triviale,
venue de la représentation : $\mathrm{Sym}^{2}\Sigma_{0}$, fibré plat de la
représentation adjointe, est $T_{X} \oplus \mathbf{1}$, alors que $T_{X}$
lui-même n'est pas plat. La page écrit d'abord « complexe », corrigé en
« réel » ; la fin de la phrase est illisible.}

\pagerange{7}{13}
\section*{Tapuscrit : le groupe de Brauer d'une fibration sur une courbe (pages 9, 11, 13, 7)}

Quatre feuillets d'un tapuscrit français, paginés 23, 17, 19 et 20 et rangés
aux pages 7, 9, 11 et 13 ; on les lit dans l'ordre de leur pagination, qui est
celui de l'argument. Ils portent des corrections de sa main --- lettres
grecques, tildes, schémas de Picard ajoutés, insertions interlinéaires, un
bloc biffé ---, et sont sans rapport avec les classes de Chern : c'est le
papier sur lequel il a écrit les feuillets 6 à 12.\footnote{La numérotation
(4.6)--(4.43), l'objet --- le groupe de Brauer d'un schéma régulier fibré en
courbes sur une courbe, le groupe de Tate--Chafarévitch --- et le style
indiquent un état dactylographié du § 4 de « Le groupe de Brauer III »
(publié dans \emph{Dix exposés sur la cohomologie des schémas}, 1968) ;
l'identification est la nôtre, et le texte n'a pas été comparé à l'imprimé.
Les définitions de $P$, $\underline{A}$, $\delta$, $\delta'$, $\delta''$,
$\Delta$, $\mathrm{III}'$ et les conditions (4.13), (4.15), (4.16) sont sur
des pages du tapuscrit qui ne sont pas dans le dossier ; on ne les devine
pas. $\mathrm{III}$ est le Ш de Tate--Chafarévitch, tapé comme un III barré.}

\emph{Cadre}, tel qu'il ressort des pages. $f : X \to Y$ est un morphisme
propre, $X$ régulier, $Y$ une courbe régulière de point générique $\eta$ ;
$P$ est le faisceau de Picard relatif $R^{1}f_{*}\mathbb{G}_{m}$ (étale) ; un
tilde désigne la localisation stricte en un point géométrique $\bar{y}$ au
dessus d'un point fermé $y$, une barre la localisation hensélienne.

\pagerange{9}{9}
\subsection*{Les faisceaux $\underline{E}$ et $\underline{F}$ (page 9, tapuscrit p.\ 17)}

La page s'ouvre sur la fin d'un corollaire ($H^{3}(Y,\mathbb{G}_{m}) \simeq
\mathbb{Q}/\mathbb{Z}$). Soit $i : \eta \to Y$ l'inclusion du point
générique et
\[
  \underline{B} = i_{*}i^{*}P, \qquad P \to \underline{B}, \qquad
  \underline{E} = \mathrm{Ker}(P \to \underline{B}), \quad
  \underline{F} = \mathrm{Coker}(P \to \underline{B}). \tag{4.8--4.10}
\]
Les faisceaux $\underline{E}$ et $\underline{F}$ sont concentrés aux points
fermés (faisceaux gratte-ciel), et leurs fibres géométriques en $\bar{y}$
sont le noyau et le conoyau de $\mathrm{Pic}(\widetilde{X}) \to
\mathrm{Pic}_{\widetilde{X}_{\tilde{\eta}}/\tilde{\eta}}(\tilde{\eta})$. Cette
flèche se factorise par $\mathrm{Pic}(\widetilde{X}_{\tilde{\eta}})$, et
$\mathrm{Pic}(\widetilde{X}) \to \mathrm{Pic}(\widetilde{X}_{\tilde{\eta}})$
est surjectif parce que $\widetilde{X}$ est régulier (tout diviseur de la
fibre générique s'étend par adhérence). D'où
\[
  \underline{F}_{\bar{y}} = \mathrm{Pic}_{\widetilde{X}_{\tilde{\eta}}/\tilde{\eta}}(\tilde{\eta})
  / \mathrm{Pic}(\widetilde{X}_{\tilde{\eta}}), \qquad
  \underline{E}_{\bar{y}} = \mathrm{Ker}\bigl(\mathrm{Pic}(\widetilde{X}) \to
  \mathrm{Pic}(\widetilde{X}_{\tilde{\eta}})\bigr). \tag{4.11--4.12}
\]
\footnote{Le tapuscrit écrit $\mathrm{Pic}(\widetilde{X}/\widetilde{Y})$ et
$\mathrm{Pic}(\widetilde{X}_{\tilde{\eta}}/\tilde{\eta})$ pour les sections
du faisceau de Picard relatif ; la main a ajouté « $\mathrm{Pic}(\widetilde{X})
=$ » devant le premier. Dans (4.11) le numérateur est tapé
$\mathrm{Pic}(\widetilde{X}_{\eta}/\tilde{\eta})$ ; on lit
$\tilde{\eta}$ partout.} Le quotient $\underline{F}_{\bar{y}}$ s'injecte dans
$\mathrm{Br}(\tilde{\eta})$, qui est nul si $k(y)$ est parfait (le corps des
fractions d'un anneau de valuation discrète strictement hensélien à corps
résiduel parfait est $C_{1}$) et de $p(y)$-torsion en général, $p(y)$ étant
l'exposant caractéristique de $k(y)$. Donc $\underline{F}_{\bar{y}}$ est de
$p(y)$-torsion, nul si $k(y)$ est parfait, et nul aussi si $\delta_{y} = 1$,
par exemple si $\widetilde{X} \to \widetilde{Y}$ a une section.

\pagerange{11}{11}
\subsection*{Le groupe de Tate--Chafarévitch (page 11, tapuscrit p.\ 19)}

Sous l'hypothèse (4.15), l'image de $H^{1}(Y,P)$ dans
$H^{1}(Y,\underline{B})$ est le sous-groupe des éléments dont l'image dans
chacun des $H^{1}(\overline{Y}_{y},\underline{B})$, $\overline{Y}_{y}$ hensélisé
de $Y$ en $y$, est nulle --- le \emph{groupe de Tate--Chafarévitch}
$\mathrm{III}(Y,\underline{B})$ : les éléments localement triviaux
partout. On en tire la suite exacte
\[
  \cdots \to \mathrm{Pic}_{X_{\eta}/\eta}(\eta) \to \bigoplus_{y}
  H^{1}(y,\underline{E}_{y}) \to H^{1}(Y,P) \to
  \mathrm{III}(Y,\underline{B}) \to 0, \tag{4.17}
\]
qui présente $H^{1}(Y,P)$ comme extension de $\mathrm{III}(Y,\underline{B})$
par un quotient de $\bigoplus_{y} H^{1}(y,\underline{E}_{y}) =
H^{1}(Y,\underline{E})$.\footnote{Un bloc d'un tiers de page, encadré et
barré, déduisait (4.15) de conditions (4.16) par un diagramme comparant les
suites du type (4.5) pour $\bar{f} : \bar{X} \to \bar{Y}$ et pour $f_{y} :
X_{y} \to y$ ; la main y avait encore corrigé (4.16) en
$\mathrm{Ker}(H^{3}(y,\mathbb{G}_{m}) \to H^{3}(X_{y},\mathbb{G}_{m}))$.
Le tapuscrit écrit $\bar{Y}$ pour le hensélisé, sans indice ; il annonce
« plus bas » la relation avec la définition originelle de Tate et
Chafarévitch, qui n'est pas dans le dossier. En marge, « (P) ».}

\emph{La structure de $\underline{E}_{y}$.} $\underline{E}_{y}$ est nul si
la fibre $X_{y}$ est géométriquement intègre ; si $X_{\eta}$ est lisse sur
$\eta$, le support de $\underline{E}$ est donc contenu dans l'ensemble fini
des $y$ dont la fibre n'est pas géométriquement intègre (le tapuscrit ajoute
« en fait, égal à »). On écrit l'égalité de diviseurs sur $X$
\[
  X_{y} = \textstyle\sum_{i} a_{y}^{i} C_{y}^{i}, \tag{4.18}
\]
les $C_{y}^{i}$ étant les composantes irréductibles de $X_{y}$ et les
$a_{y}^{i}$ leurs multiplicités.

\pagerange{13}{13}
\subsection*{L'annulation de $H^{1}(y,\underline{E}_{y})$ (page 13, tapuscrit p.\ 20)}

Soit $\mu_{y}^{i}$ et $\nu_{y}^{i}$ les multiplicités radicielle et séparable
de $C_{y}^{i}$ sur $k(y)$ ; $\mu_{y}^{i} = 1$ si $k(y)$ est parfait. Soit
$k(y)^{i}$ la fermeture séparable de $k(y)$ dans le corps des fonctions de
$C_{y}^{i}$, extension de degré $\nu_{y}^{i}$, et
\[
  Z_{y}^{i} = \mathrm{Spec}\,k(y)^{i}, \qquad Z_{y} = \textstyle\coprod_{i}
  Z_{y}^{i}, \qquad p_{y} : Z_{y} \to y. \tag{4.20}
\]
Les composantes géométriques de $C_{y}^{i}$ sont alors indexées, de façon
compatible à Galois, par les $k(y)$-plongements de $k(y)^{i}$ dans une
clôture séparable.\footnote{Le tapuscrit corrigé définit $k(y)^{i}$ comme « la
plus petite extension galoisienne de $k(y)$ telle que les composantes
irréductibles de $X_{y}^{i} \otimes_{k(y)} k(y)^{i}$ soient géométriquement
irréductibles », et lui donne le degré $\nu_{y}^{i}$. Les deux exigences ne
sont compatibles que si la fermeture séparable de $k(y)$ dans
$k(C_{y}^{i})$ est galoisienne --- ce qui est toujours le cas pour $k(y)$
fini ; sinon la plus petite extension galoisienne est sa clôture
galoisienne, de degré plus grand, et (4.21) demande la fermeture séparable
elle-même. C'est elle que donnait la première rédaction, biffée :
$Z_{y}^{i} = \mathrm{Spec}(H^{0}(X_{y}^{i},\mathcal{O}))$, au degré
inséparable près. On la garde. Le renvoi est à EGA IV 4.7 ; les numéros
(4.19) à (4.24) sont corrigés à la main sur (4.18) à (4.23).} Le faisceau
$\underline{E}_{y}$ sur $y$ s'insère dans la suite exacte
\[
  0 \to \mathbb{Z}_{y} \to p_{y*}\mathbb{Z}_{Z_{y}} \to \underline{E}_{y}
  \to 0, \tag{4.21}
\]
où la première flèche est l'adjointe de $p_{y}^{*}\mathbb{Z}_{y} =
\mathbb{Z}_{Z_{y}} \to \mathbb{Z}_{Z_{y}}$, multiplication par $a_{y}^{i}$
sur $Z_{y}^{i}$ : géométriquement, $\underline{E}_{\bar{y}}$ est le groupe
libre sur les composantes de la fibre modulo la fibre elle-même. Comme
$H^{1}$ d'un corps à coefficients dans $\mathbb{Z}$ est nul, la suite
de cohomologie donne
\[
  H^{1}(y,\underline{E}_{y}) = \mathrm{Ker}\Bigl(H^{2}(y,\mathbb{Z}) \to
  \textstyle\prod_{i} H^{2}(Z_{y}^{i},\mathbb{Z})\Bigr)
  = \mathrm{Ker}\Bigl(H^{1}(y,\mathbb{Q}/\mathbb{Z}) \to
  \textstyle\prod_{i} H^{1}(Z_{y}^{i},\mathbb{Q}/\mathbb{Z})\Bigr), \tag{4.23}
\]
la flèche d'indice $i$ étant la restriction multipliée par $a_{y}^{i}$. La
corestriction composée avec la restriction étant la multiplication par
$\nu_{y}^{i}$, chacun de ces noyaux est annulé par $a_{y}^{i}\nu_{y}^{i}$, et
$H^{1}(y,\underline{E}_{y})$ est annulé par
\[
  d_{y} = \mathrm{pgcd}_{i}\,(a_{y}^{i}\nu_{y}^{i}). \tag{4.24}
\]
Le feuillet s'arrête sur le cas où $k(y)$ est fini, ou plus généralement de
groupe de Galois absolu $\widehat{\mathbb{Z}}$ : « alors on trouve
\ldots ».\footnote{Le tapuscrit dit « pseudo-fini au sens de Serre » ; le
terme de Serre (\emph{Corps locaux}, 1962) est « quasi-fini », et
« pseudo-fini » désigne aujourd'hui la notion d'Ax (1968), plus forte. La
suite n'est pas dans le dossier. Le calcul se termine sans peine, et c'est
nous qui l'ajoutons : $H^{1}(y,\mathbb{Q}/\mathbb{Z}) = \mathbb{Q}/\mathbb{Z}$,
la restriction à une extension de degré $\nu$ est la multiplication par
$\nu$, donc $H^{1}(y,\underline{E}_{y}) \simeq \mathbb{Z}/d_{y}\mathbb{Z}$.}

\pagerange{7}{7}
\subsection*{Deux cas particuliers (page 7, tapuscrit p.\ 23)}

La page s'ouvre sur la fin d'une suite exacte
\[
  0 \to \mathbb{Z}/\delta''\mathbb{Z} \to \mathrm{III}'(Y,\underline{A}) \to
  \mathrm{III}'(Y,\underline{B}) \to 0 . \tag{4.38}
\]

\emph{(4.6) : $Y$ est une courbe algébrique sur un corps algébriquement
clos.} La condition (4.13) est automatique et l'on a
\[
  \mathrm{Br}(X) \xrightarrow{\ \sim\ } H^{1}(Y,\underline{B}) ; \tag{4.39}
\]
si de plus $X$ n'a pas de fibre multiple ($\delta_{y} = 1$ pour tout $y$),
on a $\delta = \delta' = \delta''$ et
\[
  0 \to \mathbb{Z}/\delta\mathbb{Z} \to H^{1}(Y,\underline{A}) \to
  H^{1}(Y,\underline{B}) \to 0 . \tag{4.40}
\]

\emph{(4.7) : $Y$ est une courbe propre sur un corps fini $k$.} Les conditions
(4.13) et (4.16) sont automatiques, la première parce que $k$ est parfait,
la seconde parce que $k(y)$, fini, est de dimension cohomologique $\leq 1$,
d'où $H^{3}(k(y),\mathbb{G}_{m}) = 0$ ; on a $\delta = \delta'$, et de même
pour les invariants hensélisés et strictement hensélisés. Comme
$\mathrm{Br}(Y) = 0$ (théorie du corps de classes),
\[
  0 \to \mathrm{Br}(X) \to \mathrm{III}(Y,\underline{B}) \to
  \mathbb{Z}/\Delta\mathbb{Z} \to 0, \qquad \Delta \mid \delta . \tag{4.41}
\]
Si le morphisme déduit de $f$ par extension à la clôture algébrique de $k$
n'a pas de fibre multiple, on a
\[
  0 \to \mathbb{Z}/\delta''\mathbb{Z} \to \mathrm{III}(Y,\underline{A}) \to
  \mathrm{III}(Y,\underline{B}) \to 0, \tag{4.42}
\]
cas particulier de (4.37), compte tenu de l'égalité $\bar{\delta}_{y} =
\delta_{y}$ (4.43) pour tout point fermé $y$, qui donne ici
$\bar{\delta}_{y} = 1$, donc $\mathrm{III}' = \mathrm{III}$. L'égalité (4.43)
est une conséquence du théorème de Lang, $H^{1}(k(y), G) = 0$ pour un groupe
algébrique connexe $G$ sur un corps fini, appliqué à $G =
\underline{\mathrm{Pic}}^{0}_{X_{y}/y}$ ; le tapuscrit la laisse « au lecteur
comme un exercice plaisant et délectable ».\footnote{La référence « [24]
$H^{1}(k,G) = 0$ » et « ainsi que (4.16), puisque $k(y)$ est de dimension
$\leq 1$ \ldots » sont des insertions de sa main. Les deux cas particuliers
étaient d'abord numérotés « 1 », etc. (4.41) est la forme que prend, pour
une surface fibrée sur une courbe sur un corps fini, la comparaison du groupe
de Brauer de la surface et du groupe de Tate--Chafarévitch de la jacobienne de
la fibre générique, qui est au cœur de la conjecture d'Artin--Tate ; le nom
n'est pas sur ces feuillets.}

\pagerange{15}{15}
\section*{La cohomologie réelle de $B_{G}$ (page 15)}

Soit $G$ un groupe de Lie compact connexe. L'espace classifiant est la
réalisation de la variété simpliciale $[p] \mapsto G^{p+1}/G$ (quotient par
l'action diagonale, $G^{p+1}/G \simeq G^{p}$), et sa cohomologie réelle est
celle du bicomplexe des formes différentielles
\[
  C_{(0)}^{p,q} = \Omega^{q}(G^{p+1}/G),
\]
de différentielles la différentielle de de Rham (en $q$) et la somme alternée
des faces (en $p$). On le remplace successivement par deux sous-bicomplexes de
formes invariantes, $C_{(1)} \subset C_{(0)}$ puis $C_{(2)}$, chacun ayant, à
$p$ fixé, la même cohomologie que le précédent --- c'est le terme $E_{1}$ de
la suite spectrale de la filtration par $p$, calculé par moyenne sur le groupe
compact. Dans le dernier, $C_{(2)}^{p,q} = \omega_{\mathrm{inv}}^{q}
(G^{p+1}/G)$, les formes sont fermées : la différentielle de de Rham y est
\emph{nulle}. D'où des isomorphismes canoniques
\[
  H^{*}(B_{G},\mathbb{R}) \simeq H^{*}(C_{(1)}) \simeq
  \bigoplus_{p+q = *} H^{p}\bigl(\nu \mapsto \omega_{\mathrm{inv}}^{q}(G^{\nu+1}/G)\bigr)
  = \bigoplus_{p+q = *} E_{2}^{p,q},
\]
et la suite spectrale dégénère : $d_{r} = 0$ pour $r \geq 2$.\footnote{La page
distingue $\Omega$, $\omega$ et $\omega_{\mathrm{inv}}$ sans les définir, et
dessine une inclusion de $C_{(1)}$ dans $C_{(0)}$ et une flèche de $C_{(1)}$
vers $C_{(2)}$ ; une note marginale, en partie illisible, parle de
« cohomologie pour $\nu$ fixe » et de $E_{1}$. Lire $\omega_{\mathrm{inv}}$
comme les formes bi-invariantes sur $G^{\nu} \simeq G^{\nu+1}/G$, qui sont
fermées et représentent la cohomologie de de Rham (Cartan, Chevalley--Eilenberg),
est notre interprétation ; c'est elle qui rend juste l'affirmation que la
différentielle en $q$ est nulle. Le complexe en $\nu$ est alors la
construction cobar de l'algèbre de Hopf $H^{*}(G,\mathbb{R})$. La page
affirme la dégénérescence sans la démontrer. Le bicomplexe des formes sur la
variété simpliciale $G^{\bullet+1}/G$ est aujourd'hui associé aux noms de
Bott, Shulman et Stasheff (1973) et de Dupont (1976), bien postérieurs à la
date que l'inventaire donne au dossier.}

Soit $T \subset \omega_{\mathrm{inv}}^{*}(G)$ l'espace des éléments
transgressifs --- de degrés impairs, ce sont les éléments primitifs de
l'algèbre de Hopf $H^{*}(G,\mathbb{R}) = \Lambda^{*}(T)$ --- et
\[
  \varphi : T \to E_{2}^{1,*} = \omega_{\mathrm{inv}}^{*}(G \times G/G)
  \subset \omega_{\mathrm{inv}}^{*}(G \times G) = \omega_{\mathrm{inv}}^{*}(G)
  \otimes \omega_{\mathrm{inv}}^{*}(G), \qquad
  \varphi(t) = t \otimes 1 - 1 \otimes t ,
\]
qui place un élément de degré $2k - 1$ en degré total $2k$.

\textbf{Théorème.} \emph{L'homomorphisme $\mathrm{Sym}^{*}(T) \to
H^{*}(B_{G},\mathbb{R})$ déduit de $\varphi$ est un isomorphisme.}

C'est le théorème de Borel--Cartan : $H^{*}(B_{G},\mathbb{R})$ est une algèbre
de polynômes, sur des générateurs de degrés $2k_{j}$ si $H^{*}(G,\mathbb{R})$
est extérieure sur des générateurs de degrés $2k_{j} - 1$. Le calcul de la
page en donne une démonstration par la dualité de Koszul : la cohomologie
de la construction cobar d'une algèbre extérieure $\Lambda(T)$ est l'algèbre
symétrique sur $T$ décalé d'un cran.\footnote{Le signe de $\varphi$ dépend de
l'identification de $G \times G/G$ à $G$ ($(a,b) \mapsto ab^{-1}$ donne celui
de la page). Les noms --- Borel, Cartan, Koszul --- ne sont pas sur la page,
qui écrit seulement « Th. » ; une parenthèse incertaine, « OK \ldots{} des
degrés impairs », suit la définition de $T$.}

\pagerange{16}{16}
\section*{Une lettre dactylographiée : les catégories $C(X,A)$ (page 16)}

Un feuillet de lettre dactylographiée, sans en-tête ni signature, barré de
trois traits obliques ; le destinataire est tutoyé, une lettre antérieure « du
26.3 », que le destinataire vient de renvoyer en photocopie, est citée, et une
insertion est de la main de Grothendieck.\footnote{Ni l'auteur, ni le
destinataire, ni la date ne sont sur ce feuillet ; on n'en propose pas. La
lettre parle à la première personne de « ma lettre du 26.3 » et des
« hypothèses de finitude que j'ai dites ».}

À un espace $X$ et un anneau $A$ est associée une catégorie triangulée
$C(X,A)$ --- vraisemblablement une catégorie dérivée de faisceaux localement
constants de $A$-modules, sous des conditions de finitude posées plus haut
dans la lettre.\footnote{La définition précède ce feuillet. L'interprétation
par les systèmes locaux est la nôtre ; c'est elle que suggère la mention des
« $A$-modules tordus de type fini ».} Elle détermine l'ensemble des
composantes connexes de $X$, le groupe fondamental de chacune et leur
cohomologie à coefficients dans les $A$-modules tordus de type fini. Pour $A =
\mathbb{Z}$, un morphisme $f : X \to Y$ est une équivalence d'homotopie si et
seulement si $f^{*} : C(Y,\mathbb{Z}) \to C(X,\mathbb{Z})$ est une
équivalence.\footnote{La lettre attribue ce critère, prudemment (« il semble
bien »), à Hurewicz « ou une variante mise au point par Artin--Mazur ». Pour
des CW-complexes, c'est le théorème de Whitehead à coefficients locaux : une
application qui induit des bijections sur les $\pi_{0}$ et les $\pi_{1}$ et des
isomorphismes en homologie à coefficients dans tous les systèmes locaux est
une équivalence d'homotopie. Que $f^{*}$ soit une équivalence force ces
conditions (la restriction le long de $\mathbb{Z}[\pi_{1}X] \to
\mathbb{Z}[\pi_{1}Y]$ n'est une équivalence que si cet homomorphisme est un
isomorphisme).}

La lettre demande alors si l'on peut reconstituer $X$ à partir de
$C(X,\mathbb{Z})$, du moins pour des $X$ de type fini en chaque dimension, et
si tout foncteur exact $C(Y,\mathbb{Z}) \to C(X,\mathbb{Z})$ provient, à
isomorphisme près, d'un unique morphisme $X \to Y$.\footnote{Telle qu'elle est
posée, la seconde question a une réponse négative : le foncteur nul et le
décalage $[1]$ sont exacts et ne proviennent d'aucun morphisme. Il faut au
moins demander que le foncteur respecte le produit tensoriel et l'objet
unité. Sous de telles formes, la question de reconstituer un type d'homotopie
à partir de ses coefficients a été reprise bien plus tard, par lui (\emph{À la
poursuite des champs}, 1983) et par d'autres (types d'homotopie
schématiques de Toën, théorème de Mandell sur les cochaînes entières) ; ces
rapprochements sont les nôtres.} Même si c'est trop optimiste, ces catégories
lui paraissent des invariants « rupinants », très bonnes approximations du
type d'homotopie de $X$, et il faudrait traduire en leurs termes les
constructions familières aux homotopistes : groupes d'homotopie supérieurs,
espaces de lacets, et en sens inverse espaces classifiants, suspensions,
produits, squelettes, images des groupes d'homotopie supérieurs. La question,
écrit-il, est « manifestement très liée » à celle de la lettre du 26.3, et
« quelqu'un devrait bien un jour regarder ces choses systématiquement ».

\pagerange{17}{17}
\section*{Groupes de Lie sans $p$-torsion (page 17)}

Soit $G$ un groupe de Lie complexe réductif connexe --- ou, ce qui revient au
même pour la cohomologie, son sous-groupe compact maximal ---, $B$ un
sous-groupe de Borel, $\hat{B}$ son groupe des caractères, de rang $r$, et
$p$ un nombre premier.\footnote{« compact » est biffé et remplacé par une
insertion lue « complexe réductif », lecture incertaine. $G/B$ est la variété
des drapeaux, difféomorphe à $K/T$ pour $K$ compact maximal et $T$ tore
maximal ; $B_{B} \simeq B_{T}$. La page écrit $\mathbb{B}_{B}$,
$\mathbb{B}_{G}$ pour les espaces classifiants dans la ligne sur Leray.}
Considérons les conditions :

\begin{itemize}
\item[(i)] $H^{*}(G,\mathbb{Z})$ est sans $p$-torsion ;
\item[(ii)] $H^{*}(B_{G},\mathbb{Z})$ est sans $p$-torsion ;
\item[(iii)] l'image de $\mathrm{Sym}_{\mathbb{Z}}^{*}(\hat{B}) \to
H^{*}(G/B,\mathbb{Z})$ est d'indice premier à $p$ ; de façon équivalente,
$\mathrm{Sym}_{\mathbb{F}_{p}}(\hat{B} \otimes \mathbb{F}_{p}) \to
H^{*}(G/B,\mathbb{F}_{p})$ est surjectif ;
\item[(i bis)] $H^{*}(G,\mathbb{F}_{p}) \simeq \Lambda^{*}V$, où $V$ est un
$\mathbb{F}_{p}$-espace de dimension $r$ formé d'éléments primitifs ;
\item[(ii bis)] $H^{*}(B_{G},\mathbb{F}_{p})$ est une algèbre de polynômes
sur $\mathbb{F}_{p}$ à $r$ générateurs.
\end{itemize}

Les implications notées :
\begin{itemize}
\item[--] (i) $\Rightarrow$ (i bis), par le théorème de Hopf (sous sa forme
modulo $p$, due à Borel) ;
\item[--] (i bis) $\Rightarrow$ (ii bis), par la suite spectrale de la
fibration universelle, ou le théorème de transgression de Borel ;
\item[--] (ii bis) $\Leftrightarrow$ (iii) : la suite spectrale de Leray de
la fibration $G/B \to B_{B} \to B_{G}$ est triviale ;
\item[--] (ii bis) $\Rightarrow$ (ii), par égalité des nombres de Betti à
coefficients dans $\mathbb{Q}$ et dans $\mathbb{F}_{p}$.
\end{itemize}

\begin{tikzcd}
(\mathrm{i}) \arrow[r, Rightarrow] \arrow[d, Rightarrow] &
(\mathrm{i\ bis}) \arrow[dl, Rightarrow] \\
(\mathrm{ii\ bis}) \arrow[r, leftrightarrow] \arrow[d, Rightarrow] &
(\mathrm{iii}) \\
(\mathrm{ii}) &
\end{tikzcd}
\footnote{La page renvoie à des « notes de Borel ». La dernière implication
demande que les générateurs de $H^{*}(B_{G},\mathbb{F}_{p})$ aient les degrés
des générateurs rationnels --- ce que donne la transgression, ou la
comparaison avec $B_{T}$ ; sous cette précision, les séries de Poincaré mod
$p$ et rationnelle coïncident, ce qui exclut la $p$-torsion par le théorème
des coefficients universels. Les cinq conditions sont en fait équivalentes
(Borel) ; la page n'écrit que les implications du diagramme. La flèche
horizontale du bas est dessinée à deux pointes.}

\pagerange{18}{18}
\section*{EGA IV, feuillet IV-1102 annoté (page 18)}

Un feuillet du tapuscrit d'EGA IV, numéroté IV-1102, portant des corrections
de sa main ; le passage est celui du cycle associé à un diviseur positif,
(21.6.5)--(21.6.6).\footnote{La numérotation est celle du § 21.6 d'EGA IV,
quatrième partie (1967) ; on n'a pas comparé le feuillet au texte imprimé.}
Les corrections déplacent l'hypothèse noethérienne en tête de l'alinéa
(« $X$ étant supposé noethérien », ajouté ; « Si $X$ est noethérien », biffé
dans l'énoncé de la proposition) et ajoutent les renvois.

Soit $D$ un diviseur \emph{positif} sur $X$ : l'Idéal $I_{X}(D)$ est un
Idéal inversible de $\mathcal{O}_{X}$, et définit un sous-schéma fermé
$Y(D)$. Localement $I_{X}(D) = \mathcal{O}_{U}\,t$ avec $t$ régulier
(21.2.7) : l'immersion $Y(D) \to X$ est régulière (19.1.4) de codimension 1
en tout point. Inversement, un sous-schéma fermé régulièrement immergé de
codimension 1 en tout point est $Y(D)$ pour un unique diviseur positif $D$.
\footnote{En marge de cette réciproque : « devrait figurer dans 21.2.
\ldots{} (sans noeth.) », la dernière lecture incertaine ; elle ne demande en
effet pas d'hypothèse noethérienne.}
Pour $X$ noethérien, $Y(D)$ n'a qu'un nombre fini de points maximaux, qui
sont de codimension 1 dans $X$ (théorème de l'idéal principal) et où
$\mathcal{O}_{Y(D),x}$ est artinien ; en tout autre point $x \in X^{(1)}$,
$\mathcal{O}_{Y(D),x} = 0$. On pose
\[
  \mathrm{cyc}(D) = \sum_{x \in X^{(1)}} \mathrm{long}(\mathcal{O}_{Y(D),x})
  \cdot \overline{\{x\}} \in Z^{1}(X). \tag{21.6.5.1}
\]

\textbf{Proposition (21.6.6).} \emph{$D \mapsto \mathrm{cyc}(D)$ est un
homomorphisme du monoïde $\mathrm{Div}^{+}(X)$ dans $Z^{1}(X)$.}

Comme $I_{X}(D + D') = I_{X}(D)I_{X}(D')$ (21.2.5), il suffit de voir que,
pour $A = \mathcal{O}_{X,x}$, $x \in X^{(1)}$, et $t$, $t'$ réguliers dans $A$,
$\mathrm{long}(A/tt'A) = \mathrm{long}(A/tA) + \mathrm{long}(A/t'A)$ ; or
$tA/tt'A \simeq A/t'A$ puisque $t$ est régulier, et la suite $0 \to tA/tt'A
\to A/tt'A \to A/tA \to 0$ conclut. Le feuillet s'arrête sur la phrase
suivante (« pour tout ouvert $U \subset X$, on \ldots »).\footnote{Une
seconde note marginale, « réf ! », est placée à hauteur de la définition de
$\mathrm{cyc}(D)$.}

\end{document}
